\documentclass[sigconf]{acmart}

%% Using the option "anonymous" will produce "Anonymous Author(s)"
%% for double-anonymized review. Upon acceptance, remove this option
%% and set the authors' names and affiliations using the corresponding
%% commands.

%% CBSoft uses an ACM-like format without the ACM Reference block.
%% The following warning from acmart can be safely ignored:
%% "ACM reference format is mandatory for papers over one page."
\settopmatter{printacmref=false}
\renewcommand\footnotetextcopyrightpermission[1]{} % removes footnote with conference information in first column

%% CBSoft uses an ACM-like format without a copyright note.
\setcopyright{none}

%% These commands are for a PROCEEDINGS abstract or paper.
\acmConference[SBES 2026]{40th Brazilian Symposium on Software Engineering}{September 8–11, 2026}{São Paulo, SP, Brazil}

% CBSoft uses an ACM-like format without CCS Concepts.
% The following warning from acmart can be safely ignored:
% "CCS concepts are mandatory for papers over two pages."

%% \BibTeX command to typeset BibTeX logo in the docs
\AtBeginDocument{
    \providecommand\BibTeX{{Bib\TeX}}
}

%% Only for papers written in PORTUGUESE: 
%% comment out the following three lines
\renewcommand\abstractname{Resumo}
\renewcommand\keywordsname{Palavras-chave}
\renewcommand\refname{Referências}

%% end of the preamble, start of the body of the document source.


% ------------------------------------------------------------
% Authors: please make sure to read the checklist comments
% provided at the end of this document before submission.
% ------------------------------------------------------------
\begin{document}

%% The "title" command has an optional parameter,
%% allowing the author to define a "short title" to be used on the page 
%% headers.
\title{The Name of the Title Is Hope}

%% The "author" command and its associated commands are used to define
%% the authors and their affiliations.
%% Of note is the shared affiliation of the first two authors, and the
%% "authornote" and "authornotemark" commands
%% used to denote shared contribution to the research.
\author{FirstName Surname1stAutor}
\affiliation{
  \institution{DepartmentName, Institution/UniversityName}
  \city{City}
  \country{Country}
}
\email{email@email.com}

\author{FirstName Surname}
\affiliation{
  \institution{DepartmentName, Institution/UniversityName}
  \city{City}
  \country{Country}
}
\email{email@email.com}

\author{FirstName Surname}
\affiliation{
  \institution{DepartmentName, Institution/UniversityName}
  \city{City}
  \country{Country}
}
\email{email@email.com}

\author{FirstName Surname}
\affiliation{
  \institution{DepartmentName, Institution/UniversityName}
  \city{City}
  \country{Country}
}
\email{email@email.com}

\author{FirstName Surname}
\affiliation{
  \institution{DepartmentName, Institution/UniversityName}
  \city{City}
  \country{Country}
}
\email{email@email.com}

\author{FirstName Surname}
\affiliation{
  \institution{DepartmentName, Institution/UniversityName}
  \city{City}
  \country{Country}
}
\email{email@email.com}

\author{FirstName Surname}
\affiliation{
  \institution{DepartmentName, Institution/UniversityName}
  \city{City}
  \country{Country}
}
\email{email@email.com}

\author{FirstName Surname}
\affiliation{
  \institution{DepartmentName, Institution/UniversityName}
  \city{City}
  \country{Country}
}
\email{email@email.com}

\author{FirstName Surname}
\affiliation{
  \institution{DepartmentName, Institution/UniversityName}
  \city{City}
  \country{Country}
}
\email{email@email.com}

%% By default, the full list of authors will be used on the page
%% headers. This list is often too long and will overlap
%% other information printed in the page headers. 
%% This command allows the author to define a more concise list
%% of authors' names for this purpose.
\renewcommand{\shortauthors}{Surname1stAuthor et al.}
\renewcommand{\shorttitle}{The Name of the Title Is Hope}

%% The abstract is a short summary of the work the paper presents.
\begin{abstract}
A evolução contínua de sistemas de software frequentemente leva ao acúmulo de anti-patterns arquiteturais, comprometendo atributos importantes de qualidade interna, como acoplamento, coesão e testabilidade. Embora a literatura apresente diversas recomendações para mitigação desses problemas, ainda existe a necessidade de evidências empíricas sobre os impactos da remoção sistemática de anti-patterns em sistemas reais.

Este trabalho apresenta um estudo de caso conduzido em uma aplicação móvel desenvolvida em React Native que apresentava problemas arquiteturais relacionados a dependências excessivas entre módulos, responsabilidades dispersas em componentes de renderização de interface, ausência de mecanismos de abstração e inversão de dependência. Para tratar esses problemas, foi realizado um processo de refatoração arquitetural baseado em princípios de Engenharia de Software, incluindo a separação de responsabilidades, a introdução de camadas de abstração e a reorganização da lógica de negócio.

O impacto das modificações foi avaliado por meio da comparação de métricas de qualidade coletadas antes e após a refatoração, considerando os atributos de acoplamento, coesão e testabilidade. Os resultados obtidos indicam que a remoção sistemática dos anti-patterns contribuiu para a melhoria da qualidade interna do sistema, reduzindo dependências indevidas entre componentes, aumentando a especialização dos módulos e facilitando a construção de testes automatizados.

Os resultados reforçam a importância da refatoração arquitetural como estratégia para evolução de sistemas existentes e fornecem evidências sobre os benefícios da aplicação de práticas de Engenharia de Software na melhoria da manutenibilidade de aplicações móveis.
\end{abstract}

%% Keywords. The author(s) should pick words that accurately describe
%% the presented work. Separate the keywords with commas.
\keywords{Do, Not, Use, This, Code, Put, the, Correct, Terms, for, Your, Paper}

%% This command processes the author, affiliation, and title
%% information and builds the first part of the formatted document.
%% "frenchspacing" avoids an additional space after a period at the end of a sentence.
\frenchspacing
\maketitle

\section{Introduction}

A evolução contínua de sistemas de software frequentemente resulta no acúmulo de dívida técnica e na introdução de decisões arquiteturais que comprometem a qualidade interna do sistema. À medida que novas funcionalidades são incorporadas e prazos de entrega se tornam mais restritivos, é comum que soluções inicialmente adequadas passem a apresentar problemas estruturais que dificultam a manutenção e a evolução do software.

Entre esses problemas, destacam-se os anti-patterns arquiteturais, definidos como soluções recorrentes para problemas de projeto que, apesar de aparentemente eficazes em um primeiro momento, produzem consequências negativas no longo prazo. Exemplos comuns incluem o acoplamento excessivo entre componentes, a concentração de múltiplas responsabilidades em um único módulo e a mistura de regras de negócio com detalhes de infraestrutura ou interface de usuário.

A presença desses anti-patterns pode impactar diretamente atributos de qualidade interna, como acoplamento, coesão e testabilidade. Sistemas excessivamente acoplados tendem a apresentar maior dificuldade de manutenção e evolução, enquanto módulos com baixa coesão frequentemente acumulam responsabilidades distintas e tornam o comportamento do sistema mais difícil de compreender. Da mesma forma, a baixa testabilidade reduz a capacidade de validar o comportamento do software de forma automatizada, aumentando os riscos associados à introdução de modificações.

Embora a literatura apresente diversas práticas e princípios voltados à melhoria da qualidade arquitetural, ainda existe a necessidade de evidências empíricas sobre os impactos concretos da remoção sistemática de anti-patterns em sistemas reais. Nesse contexto, este trabalho apresenta um estudo de caso conduzido na aplicação móvel XXXXXX, desenvolvida originalmente como parte de um Trabalho de Conclusão de Curso (TCC). O sistema tem como objetivo auxiliar usuários na redução do tempo de uso de smartphones por meio de mecanismos de gamificação e desafios coletivos. Durante sua evolução, foram identificados diversos problemas arquiteturais que comprometiam atributos importantes de qualidade interna, tornando o sistema um cenário adequado para investigar os efeitos de um processo estruturado de refatoração arquitetural.

O objetivo deste trabalho é avaliar o impacto da remoção sistemática de anti-patterns arquiteturais sobre atributos de qualidade interna de software. Para isso, são analisadas métricas relacionadas a acoplamento, coesão e testabilidade antes e após a realização das intervenções arquiteturais. Os resultados obtidos permitem discutir os benefícios e limitações da aplicação de práticas de Engenharia de Software voltadas à melhoria da estrutura interna de aplicações móveis.



\section{Fundamentação Teórica}

\subsection{Arquitetura de Software e Qualidade Interna}

A arquitetura de software corresponde à estrutura fundamental de um sistema, compreendendo seus componentes, as relações entre eles e os princípios que orientam sua organização e evolução. Segundo Bass, Clements e Kazman~\cite{Bass2012}, a arquitetura representa um conjunto de estruturas necessárias para raciocinar sobre o sistema, incluindo elementos de software, propriedades visíveis externamente e os relacionamentos estabelecidos entre esses elementos.

A definição de uma arquitetura adequada desempenha papel fundamental na qualidade de sistemas de software, influenciando atributos como manutenibilidade, escalabilidade, testabilidade e capacidade de evolução. Decisões arquiteturais tomadas durante o desenvolvimento afetam diretamente a forma como novas funcionalidades são incorporadas e como modificações futuras podem ser realizadas.

Entretanto, à medida que sistemas evoluem, é comum que restrições de prazo, mudanças de requisitos e decisões de implementação introduzam degradações na estrutura arquitetural originalmente planejada. Esse processo pode resultar no surgimento de dependências inadequadas, responsabilidades mal distribuídas e outras características que dificultam a manutenção e a evolução do software.

Nesse contexto, a arquitetura exerce papel central na sustentabilidade de longo prazo dos sistemas, tornando necessária a adoção de práticas de monitoramento e refatoração capazes de preservar atributos de qualidade interna ao longo de sua evolução.


\subsection{Anti-Patterns Arquiteturais}

Anti-patterns são soluções recorrentes para problemas recorrentes que produzem consequências negativas para o sistema. O conceito foi formalizado por Brown et al.~\cite{Brown1998}, que os definiram como uma forma literária que descreve uma solução comumente utilizada para um problema recorrente e que gera consequências negativas. Ao contrário dos design patterns, que representam soluções estruturalmente eficazes, anti-patterns descrevem soluções aparentemente funcionais que, com o tempo, comprometem atributos de qualidade interna do sistema.

Importa distinguir anti-patterns arquiteturais de code smells. Code smells, descritos por Fowler~\cite{Fowler2018}, são indícios de problemas locais no código-fonte — como métodos excessivamente longos ou variáveis com nomes genéricos — que podem ser corrigidos por refatorações pontuais. Anti-patterns arquiteturais, em contraste, referem-se a problemas estruturais que comprometem a organização das responsabilidades entre módulos e camadas do sistema, exigindo intervenções de maior escopo para sua remoção.

No contexto deste trabalho, foram identificados os seguintes anti-patterns no sistema estudado:

\textbf{Business Logic in UI.} Ocorre quando regras de negócio são implementadas diretamente em componentes de interface de usuário, em vez de residirem em uma camada dedicada de aplicação ou domínio. Esse padrão viola o Princípio da Separação de Responsabilidades e dificulta a escrita de testes automatizados, uma vez que verificar a lógica de negócio exige renderizar a interface completa.

\textbf{God Object.} Consiste na existência de módulos que concentram múltiplas responsabilidades distintas, tornando-se excessivamente grandes e difíceis de compreender. A presença de um God Object viola o Princípio da Responsabilidade Única~\cite{Martin2017} e amplifica o impacto de modificações, pois alterações em uma responsabilidade tendem a afetar as demais.

\textbf{Hardcoded Dependencies.} Ocorre quando módulos dependem diretamente de implementações concretas de serviços externos, em vez de dependerem de abstrações. Esse padrão impede a substituição das dependências por implementações alternativas — como mocks em testes — e aumenta o acoplamento entre camadas do sistema.

\textbf{Singleton Não Injetável.} Refere-se ao uso de instâncias globais únicas exportadas como valores estáticos, sem suporte à injeção de dependência. Embora o padrão Singleton seja amplamente documentado, seu uso sem mecanismo de substituição impede a escrita de testes unitários isolados, comprometendo diretamente a testabilidade do sistema~\cite{Feathers2004}.


\subsection{Qualidade Interna: Acoplamento, Coesão e Testabilidade}

A qualidade de um sistema de software pode ser avaliada sob diferentes perspectivas. A ISO/IEC~25010~\cite{ISO25010} organiza os atributos de qualidade em características externas — perceptíveis pelo usuário final — e características internas — observáveis na estrutura do código-fonte. Este trabalho concentra-se em três atributos de qualidade interna: acoplamento, coesão e testabilidade.

\textbf{Acoplamento} corresponde ao grau com que um módulo depende de implementações concretas de outros módulos. Sistemas com alto acoplamento são mais difíceis de modificar e testar, pois alterações em um módulo tendem a se propagar para outros. Uma métrica formal amplamente utilizada é o \textit{Coupling Between Object classes} (CBO), proposta por Chidamber e Kemerer~\cite{Chidamber1994}, que contabiliza o número de classes às quais uma classe está diretamente acoplada.

\textbf{Coesão} refere-se ao grau com que os elementos de um módulo estão relacionados entre si e contribuem para uma única responsabilidade bem definida. Alta coesão indica que um módulo realiza uma função específica, enquanto baixa coesão indica acúmulo de responsabilidades distintas em um mesmo artefato. O Princípio da Responsabilidade Única (SRP), enunciado por Martin~\cite{Martin2017}, estabelece que um módulo deve ter apenas uma razão para mudar. A métrica \textit{Lack of Cohesion of Methods} (LCOM), também de Chidamber e Kemerer~\cite{Chidamber1994}, quantifica esse atributo.

\textbf{Testabilidade} descreve a facilidade com que um módulo pode ser verificado de forma isolada, sem necessidade de configurar dependências externas reais. Feathers~\cite{Feathers2004} argumenta que sistemas com baixa testabilidade frequentemente apresentam dependências diretas a recursos externos, ausência de injeção de dependência e uso extensivo de singletons — características que tornam inviável a escrita de testes unitários rápidos e determinísticos.

Os três atributos são inter-relacionados: alto acoplamento tende a comprometer tanto a coesão quanto a testabilidade, enquanto a introdução de abstrações contribui simultaneamente para a redução do acoplamento e para o aumento da testabilidade.


\subsection{Refatoração Arquitetural}

Refatoração é o processo de alterar a estrutura interna de um software sem modificar seu comportamento externamente observável~\cite{Fowler2018}. Fowler catalogou diversas técnicas de refatoração aplicáveis ao nível de código — como extração de método, substituição de variável temporária por consulta e decomposição de condicional. Contudo, quando o sistema acumula anti-patterns arquiteturais, refatorações pontuais são insuficientes.

Refatoração arquitetural difere da refatoração de código em escopo e impacto. Enquanto refatorações de código atuam em construções individuais, refatorações arquiteturais envolvem a reorganização de responsabilidades entre módulos e camadas do sistema. Esse tipo de intervenção é necessário quando os problemas estruturais não podem ser resolvidos sem alterar as fronteiras entre componentes.

Neste trabalho, o processo de refatoração arquitetural envolveu a aplicação dos seguintes padrões:

\textbf{Use Cases (Camada de Aplicação).} Segundo Martin~\cite{Martin2017}, use cases representam os fluxos de negócio da aplicação e devem residir em uma camada separada da interface de usuário e da infraestrutura. Extrair use cases das telas resolve o anti-pattern \textit{Business Logic in UI} e permite testar a lógica de negócio de forma independente da UI.

\textbf{Ports \& Adapters (Arquitetura Hexagonal).} Proposta por Cockburn~\cite{Cockburn2005}, a Arquitetura Hexagonal organiza o sistema em torno de um núcleo de domínio isolado por interfaces denominadas \textit{ports}. As implementações concretas das dependências externas — os \textit{adapters} — implementam essas interfaces e são fornecidas ao núcleo por injeção de dependência. Esse padrão resolve diretamente o anti-pattern \textit{Hardcoded Dependencies}.

\textbf{Entidades de Domínio Ricas.} No contexto do \textit{Domain-Driven Design} (DDD), Evans~\cite{Evans2003} defende que entidades e objetos de valor devem encapsular as regras de negócio que lhes pertencem, em vez de serem estruturas de dados passivas com lógica dispersa em serviços. Entidades de domínio ricas aumentam a coesão ao reunir dados e comportamento relacionados em um mesmo artefato.

\textbf{Organização por Feature (Vertical Slicing).} Em vez de organizar o código por tipo de artefato — todas as telas em um diretório, todos os serviços em outro —, a organização por feature agrupa em um mesmo módulo todos os artefatos pertencentes a uma mesma funcionalidade do sistema. Essa abordagem melhora a rastreabilidade das mudanças e reduz o impacto de modificações dentro de cada domínio funcional.


\section{Metodologia}

\subsection{Tipo de Pesquisa}

Este trabalho adota a estratégia de estudo de caso único de natureza exploratório-descritiva~\cite{Yin2014}. Segundo Yin~\cite{Yin2014}, estudos de caso são especialmente adequados quando se busca compreender \textit{como} e \textit{por quê} determinadas mudanças produzem efeitos em um contexto bem delimitado — o que corresponde precisamente à questão de pesquisa investigada neste trabalho. A existência de acesso irrestrito ao código-fonte nos dois estados do sistema, aliada à documentação detalhada de cada intervenção em arquivos de especificação individuais e a um histórico de commits atômicos, fornece as condições necessárias para uma análise antes-e-depois confiável, tornando o estudo de caso a estratégia mais adequada dentre as disponíveis.

A natureza exploratória da investigação reflete a ausência de estudos similares que examinem o impacto de refatorações arquiteturais combinadas — Use Cases, Ports \& Adapters e entidades de domínio ricas — em aplicações React Native com Firebase como plataforma de back-end. O componente descritivo, por sua vez, corresponde ao relato sistemático dos indicadores coletados antes e após as intervenções, com o objetivo de gerar evidências transferíveis para contextos similares.

\subsection{Objeto de Estudo}

O objeto de estudo é a aplicação móvel XXXXXX, desenvolvida em React Native com TypeScript. O sistema utiliza Firebase como plataforma de back-end, por meio dos serviços Firestore, Authentication, Storage e Cloud Messaging, e inclui um módulo nativo Android que acessa as APIs \textit{AccessibilityService} e \textit{UsageStatsManager} do sistema operacional para monitoramento e bloqueio de aplicativos. O sistema foi desenvolvido originalmente como Trabalho de Conclusão de Curso e evoluiu de forma orgânica ao longo de múltiplas iterações, sem a adoção sistemática de práticas arquiteturais formais — o que resultou no acúmulo progressivo dos anti-patterns descritos na seção anterior.

Esse histórico torna o sistema um objeto de estudo representativo de uma classe comum de aplicações: projetos que nascem com escopo restrito e crescem em complexidade funcional sem que a estrutura interna seja sistematicamente revisada. O escopo da análise compreende o código da aplicação cliente no diretório \texttt{app/}, incluindo telas, serviços, contextos de estado, hooks, entidades de domínio, ports, adapters e use cases. O módulo nativo Android e as Cloud Functions são referenciados pontualmente quando relevantes para a análise de acoplamento, mas não constituem o foco principal dos indicadores coletados.

\subsection{Procedimento de Coleta de Dados}

Os dados foram coletados em dois pontos no tempo. O baseline foi registrado em 11 de junho de 2026, imediatamente antes da execução das intervenções arquiteturais, e corresponde ao estado do sistema após um longo período de desenvolvimento sem refatoração estruturada. O snapshot pós-refatoração foi registrado em 13 de junho de 2026, após a conclusão das onze etapas do plano de melhorias. Em ambos os pontos, a coleta foi realizada por inspeção direta do código-fonte, combinando leitura manual dos arquivos com contagens automatizadas de linhas, declarações de importação e ocorrências de padrões específicos por meio de ferramentas de busca textual no repositório.

Os indicadores foram organizados em três dimensões — acoplamento, coesão e testabilidade — e operacionalizados como contagens diretas sobre o código-fonte. Optou-se por indicadores de natureza puramente observacional, sem pontuações ponderadas ou escalas compostas, pois o ecossistema TypeScript não dispõe de ferramentas de análise estática com suporte formal às métricas CBO e LCOM. Essa escolha aumenta a reprodutibilidade dos resultados: qualquer pesquisador com acesso ao repositório pode verificar os valores coletados por meio das mesmas buscas textuais.

\subsection{Indicadores de Acoplamento}
\label{sec:indicadores-acoplamento}

O acoplamento foi avaliado por dois conjuntos complementares de indicadores, correspondentes a dois níveis estruturais distintos do sistema. O primeiro conjunto mede a presença do acoplamento na camada existente, contabilizando quantos artefatos ainda dependem diretamente do Firebase SDK sem mediação por abstrações. O segundo conjunto mede o isolamento do núcleo arquitetural novo, verificando se os use cases, ports e adapters introduzidos pela refatoração mantêm o Firebase fora de suas dependências diretas. Essa distinção é necessária porque a refatoração foi incremental: parte do sistema foi migrada para o novo padrão, enquanto outra parte permanece operacional na forma anterior. A Tabela~\ref{tab:indicadores-acoplamento} apresenta os valores coletados nos dois pontos no tempo.

\begin{table}[h]
  \caption{Indicadores de acoplamento coletados nos dois pontos no tempo}
  \label{tab:indicadores-acoplamento}
  \begin{tabular}{lcc}
    \toprule
    Indicador & Baseline & Pós-refat. \\
    \midrule
    Telas importando Firebase SDK diretamente & 8 / 37 & 7 / 34 \\
    Services acoplados ao Firebase SDK & 8 / 8 & 8 / 8 \\
    Credenciais hardcoded no código-fonte & 1 & 0 \\
    Interfaces de porta (\textit{ports}) definidas & 0 & 4 \\
    Adapters concretos implementados & 0 & 4 \\
    Use cases dependentes exclusivamente de interfaces & 0 & 3 \\
    \bottomrule
  \end{tabular}
\end{table}

A aparente estabilidade nos dois primeiros indicadores — telas com Firebase direto e serviços acoplados ao SDK — reflete que as telas mais complexas e todos os serviços legados não foram migrados nesta rodada de refatoração: eles continuam funcionando sem alteração e representam a dívida técnica remanescente. No núcleo arquitetural novo, contudo, nenhum use case, entidade de domínio ou port importa o Firebase SDK diretamente. Todas as dependências externas são acessadas exclusivamente por meio das interfaces definidas em \texttt{app/ports/} e implementadas nos adapters concretos em \texttt{app/adapters/}. Essa separação representa o efeito arquitetural central da introdução do padrão Ports \& Adapters: o Firebase deixa de ser uma dependência estrutural do núcleo da aplicação e passa a ser um detalhe de infraestrutura substituível.

\subsection{Indicadores de Coesão}
\label{sec:indicadores-coesao}

A coesão foi avaliada por indicadores que capturam tanto a concentração excessiva de responsabilidades — identificada pelo tamanho e pela multiplicidade funcional dos arquivos — quanto a introdução de artefatos com responsabilidades bem delimitadas, como use cases e entidades de domínio. A Tabela~\ref{tab:indicadores-coesao} apresenta os valores coletados.

\begin{table}[h]
  \caption{Indicadores de coesão coletados nos dois pontos no tempo}
  \label{tab:indicadores-coesao}
  \begin{tabular}{lcc}
    \toprule
    Indicador & Baseline & Pós-refat. \\
    \midrule
    Arquivos $>$ 500 linhas com múltiplas responsabilidades & 3 & 2 \\
    Maior arquivo do sistema (linhas) & 2.705 & 2.696 \\
    \texttt{groupService.ts} (linhas) & 678 & 420 \\
    Total de linhas em \texttt{app/services/} & 2.596 & 2.245 \\
    Use cases criados & 0 & 3 \\
    Entidades de domínio extraídas & 0 & 4 \\
    Custom hooks de lógica de negócio & 0 & 2 \\
    \bottomrule
  \end{tabular}
\end{table}

A redução de 678 para 420 linhas no \texttt{groupService.ts} e a queda de 2.596 para 2.245 linhas no total do diretório \texttt{app/services/} refletem o movimento de lógica de negócio para artefatos mais especializados. A criação de três use cases, quatro entidades de domínio e dois hooks de lógica de negócio representa a introdução de módulos com responsabilidade única e bem delimitada, cada um orquestando ou encapsulando apenas o comportamento para o qual foi projetado. O maior arquivo do sistema — \texttt{DetalhesDoGrupoScreen.tsx}, com 2.696 linhas — e os dois God Objects remanescentes (\texttt{screenTime.ts} e \texttt{statisticsService.ts}) ilustram a dívida técnica que não pôde ser endereçada nesta rodada de refatoração, dado o escopo do trabalho, e constituem o principal alvo de uma intervenção futura.

\subsection{Indicadores de Testabilidade}
\label{sec:indicadores-testabilidade}

A testabilidade foi avaliada por indicadores que medem a capacidade de isolar unidades de código para verificação automatizada, sem necessidade de configurar Firebase real, emular um dispositivo físico Android ou instanciar os contextos globais da aplicação. A Tabela~\ref{tab:indicadores-testabilidade} apresenta os valores coletados.

\begin{table}[h]
  \caption{Indicadores de testabilidade coletados nos dois pontos no tempo}
  \label{tab:indicadores-testabilidade}
  \begin{tabular}{lcc}
    \toprule
    Indicador & Baseline & Pós-refat. \\
    \midrule
    Singletons não injetáveis nos services & 2 & 1 \\
    Use cases com injeção de dependência via interfaces & 0 & 3 \\
    Entidades de domínio sem dependências externas & 0 & 4 \\
    Ocorrências de erros tipados (\texttt{DomainError} et al.) & 0 & 47 \\
    Retornos silenciosos (\texttt{return null/false/[]}) & presente & 41 \\
    Suites de teste automatizadas existentes & 0 & 4 \\
    \bottomrule
  \end{tabular}
\end{table}

O indicador de maior impacto estrutural é a introdução de três use cases com injeção de dependência: a lógica de negócio que antes só podia ser exercida ao renderizar a tela completa — com todos os providers de contexto e mocks do Firebase configurados — pode agora ser verificada de forma isolada, fornecendo implementações de teste das interfaces \texttt{IUserRepository} e \texttt{IScreenTimeGateway}. As quatro entidades de domínio criadas são estruturas sem dependências externas, o que as torna verificáveis com testes unitários simples. A introdução de 47 ocorrências de erros tipados por meio da hierarquia \texttt{DomainError} torna os caminhos de falha explícitos e verificáveis em testes, em contraste com o comportamento anterior de retornos silenciosos — padrão que ainda persiste em 41 locais nos serviços legados, onde a lógica permanece difícil de testar de forma isolada.

\subsection{Processo de Refatoração}
\label{sec:processo-refatoracao}

A refatoração foi executada em onze etapas sequenciais, formalizadas antes do início da implementação em um documento de planejamento composto por arquivos de especificação individuais para cada melhoria. A sequência foi definida com base na relação entre o impacto esperado sobre os atributos de qualidade e o esforço de implementação estimado, e as dependências entre etapas foram respeitadas para garantir que cada melhoria pudesse ser aplicada sobre uma base estável. Cada etapa foi executada em um commit atômico no repositório, criando um histórico rastreável que permite isolar o efeito de cada intervenção individualmente.

A ordem de execução obedeceu a uma progressão do geral para o específico: a remoção de código de infraestrutura desnecessário e a padronização do tratamento de erros precederam a extração de use cases, que por sua vez precedeu a criação da camada Ports \& Adapters. Essa sequência é justificada pelo fato de que use cases com contratos de erro bem definidos são mais fáceis de isolar por interfaces — a inversão dessa ordem introduziria retrabalho. A Tabela~\ref{tab:melhorias} lista as seis intervenções de maior impacto e as dimensões de qualidade por elas afetadas.

\begin{table}[h]
  \caption{Principais intervenções e dimensões de qualidade afetadas}
  \label{tab:melhorias}
  \begin{tabular}{lc}
    \toprule
    Intervenção & Dimensões afetadas \\
    \midrule
    Padronizar tratamento de erros & Testabilidade \\
    Extrair use cases & Coesão, Testabilidade \\
    Criar camada Ports \& Adapters & Acoplamento, Testabilidade \\
    Introduzir entidades de domínio & Coesão, Testabilidade \\
    Reorganizar por feature & Coesão \\
    Mover lógica para Cloud Functions & Acoplamento, Coesão \\
    \bottomrule
  \end{tabular}
\end{table}

\section{Conclusão e Trabalhos Futuros}
<text included by the authors>

\section*{Artifact Availability}
This section must be included immediately after the Conclusion and before the Acknowledgments section. It must not be numbered and must be titled exactly as \textbf{Artifact Availability}.

If the paper provides research artifacts (such as datasets, source code, scripts, models, or experimental packages), the authors must describe what is available, where it can be accessed, and under which license or usage conditions.

If no artifacts are available, this section must still be included, explicitly stating that no artifacts are provided.

For papers written in Portuguese, this section must be titled
\textbf{Disponibilidade de Artefatos}.

\section*{Acknowledgements}
<text included by the authors>

%% The next two lines define the bibliography style to be used, and
%% the bibliography file.
\bibliographystyle{ACM-Reference-Format}
\bibliography{sample-base}

%% If your work has appendices, this is the place to put them
\appendix

\section{Appendix Example 1}

\subsection{Part One}
Lorem ipsum dolor sit amet, consectetur adipiscing elit. Morbi malesuada, quam in pulvinar varius, metus nunc fermentum urna, id sollicitudin purus odio sit amet enim. Aliquam ullamcorper eu ipsum vel mollis. Curabitur quis dictum nisl. Phasellus vel semper risus, et lacinia dolor. Integer ultricies commodo sem nec semper.

\subsection{Part Two}
Etiam commodo feugiat nisl pulvinar pellentesque. Etiam auctor sodales ligula, non varius nibh pulvinar semper. Suspendisse nec lectus non ipsum convallis congue hendrerit vitae sapien. Donec at laoreet eros. Vivamus non purus placerat, scelerisque diam eu, cursus ante. Etiam aliquam tortor auctor efficitur mattis.

\section{Appendix Example 2}
Nam id fermentum dui. Suspendisse sagittis tortor a nulla mollis, in pulvinar ex pretium. Sed interdum orci quis metus euismod, et sagittis enim maximus. Vestibulum gravida massa ut felis suscipit congue. Quisque mattis elit a risus ultrices commodo venenatis eget dui. Etiam sagittis eleifend elementum.

Nam interdum magna at lectus dignissim, ac dignissim lorem rhoncus. Maecenas eu arcu ac neque placerat aliquam. Nunc pulvinar massa et mattis lacinia.

\end{document}

% ============================================================
% CBSoft ACM-like Template – Author Checklist
% ============================================================
%
% 1. Remove CCS Concepts
%    - Remove the CCSXML block and \ccsdesc commands.
%
% 2. Remove ACM footer
%    Use the command \renewcommand\footnotetextcopyrightpermission[1]{}
%
% 3. Remove ACM Reference Format
%    Use the commands \setcopyright{none} and 
%    \settopmatter{printacmref=false}
%
% 4. Short title must not exceed column width
%    Use the command 
%    \renewcommand{\shorttitle}{The Name of the Title Is Hope}
%
% 5. Include authors' surnames in the page headers (even pages)
%    - The author names in the page header (at the top of the page)
%      must not exceed the column width
%    - This list must contain only the authors' surnames
%    - Suggestion: use "Surname et al." for more than two authors with the
%      command \renewcommand{\shortauthors}{Surname1stAuthor et al.}
%
% 6. Author format below title:
%    \author{FirstName Surname}
%    \affiliation{
%      \institution{Department, Institution}
%      \city{City}
%      \country{Country}}
%    \email{email@email.com}
%
% 7. Verify all author emails are included
%    - Check font consistency
%
% 8. All numbered section and subsection titles must follow the same
%    capitalization style: the first letter of each word in the title 
%    must be capitalized (Title Case).
%
% 9. Artifact Availability (check CFP)
%    - Not numbered -- use \section*{Artifact Availability}
%    - Placed after Conclusion and before Acknowledgments
%    - In Portuguese, use \section*{Disponibilidade de Artefatos}
%
% 10. Abstract
%     - Single paragraph
%     - Papers written in Portuguese must use "Resumo" instead of "Abstract"
%       with the command \renewcommand\abstractname{Resumo}
%
% 11. Keywords
%     - Comma-separated, first letter capitalized
%     - Papers written in Portuguese must use "Palavras-chave" instead of
%       "Keywords" with the command 
%       \renewcommand\keywordsname{Palavras-chave}
%
% 12. References
%     - Not numbered -- already enforced by the template
%     - Papers written in Portuguese must use "Referências" instead of
%       "References" with the command \renewcommand\refname{Referências}
%
% 13. Acknowledgements
%     - Optional section
%     - Not numbered -- use \section*{Acknowledgements}
%     - In Portuguese, use \section*{Agradecimentos}
%
% 14. Remove received/revised/accepted dates
%
% 15. Figure captions must appear below figures
%
% 16. Verify the symposium name in the page header
%    - Use English, even if the paper is written in Portuguese
%    - Select the correct event according to the track
%    - Use Arabic ordinal numbers, not Roman
%
%    SBES (Research, Education, IIER, Tools, Industry)
%    \acmConference[SBES 2026]{40th Brazilian Symposium on Software Engineering}{September 8--11, 2026}{São Paulo, SP, Brazil}
%
%    SBLP
%    \acmConference[SBLP 2026]{30th Brazilian Symposium on Programming Languages}{September 8--11, 2026}{São Paulo, SP, Brazil}
%
%    SBCARS
%    \acmConference[SBCARS 2026]{20th Brazilian Symposium on Software Components, Architectures, and Reuse}{September 8--11, 2026}{São Paulo, SP, Brazil}
%
%    SAST
%    \acmConference[SAST 2026]{11st Brazilian Symposium on Systematic and Automated Software Testing}{September 8--11, 2026}{São Paulo, SP, Brazil}
%
%    Workshops
%    \acmConference[Workshop Acronym 2026]{Workshop Full Name}{Workshop Date}{São Paulo, SP, Brazil}
%
%
% 17. Remove colored author names, citations, and links
%
% 18. Camera-ready only:
%     Do not use the "review" or "anonymous" modes in the 
%     \documentclass command
%
% 19. SBES Tools Track papers must include a demo video link at the end of
%     the abstract
%     Example:
%     \begin{abstract}
%     <abstract text> \newline
%     \textbf{Demo video:} <video link>
%     \end{abstract}
%
% 20. SBES Industry Track papers must not include author biographies
%
% ============================================================

\endinput


